\documentclass[11pt]{article}
\usepackage[utf8]{inputenc}
\usepackage{amsmath,amssymb,amsthm}
\usepackage[margin=1in]{geometry}

\newtheorem{lemma}{Lemma}
\newtheorem{theorem}{Theorem}

\title{Problem 655: Divisible Palindromes}
\author{}
\date{}

\begin{document}
\maketitle

\section{Problem Statement}

Count the palindromes below $10^{32}$ that are divisible by $10{,}000{,}019$.

\section{Mathematical Foundation}

Fix a digit length $d \in \{1,\dots,32\}$ and write a $d$-digit palindrome as
\[
  n = \sum_{i=0}^{h-1} x_i w_i,
\]
where $h = \lceil d/2 \rceil$ and $x_0,\dots,x_{h-1}$ are the free digits. Here
\[
  w_i =
  \begin{cases}
    10^{d-1-i} + 10^i & \text{if } i < d-1-i,\\
    10^{d-1-i} & \text{if } i = d-1-i.
  \end{cases}
\]

\begin{lemma}
A $d$-digit palindrome is uniquely determined by its first $h=\lceil d/2\rceil$ digits.
\end{lemma}

\begin{proof}
The remaining digits are forced by the symmetry relation $a_i = a_{d-1-i}$.
\end{proof}

\begin{lemma}
For fixed $d$, counting palindromes divisible by $M=10{,}000{,}019$ is equivalent to counting digit vectors $(x_0,\dots,x_{h-1})$ with $x_0 \in \{1,\dots,9\}$, $x_i \in \{0,\dots,9\}$ for $i \ge 1$, and
\[
  \sum_{i=0}^{h-1} x_i w_i \equiv 0 \pmod M.
\]
\end{lemma}

\begin{proof}
By the previous lemma each palindrome corresponds to exactly one free-digit vector, and divisibility by $M$ is exactly the displayed congruence.
\end{proof}

\begin{theorem}
The count for each length $d$ can be computed by meet-in-the-middle.
\end{theorem}

\begin{proof}
Split the free digits into left and right blocks:
\[
  \sum_{i=0}^{h-1} x_i w_i
  =
  \sum_{i \in L} x_i w_i + \sum_{i \in R} x_i w_i.
\]
For every right-half assignment, store its residue modulo $M$. Then enumerate all left-half assignments and count how many right-half residues equal the additive inverse of the left residue. Every full assignment splits uniquely into its two halves, so this counts each valid palindrome exactly once.
\end{proof}

\section{Algorithm}

For each length $d=1,\dots,32$:
\begin{enumerate}
  \item Compute the weights $w_i \bmod M$.
  \item Split the $h$ free positions into two parts, with the right part of size at most $8$.
  \item Enumerate all right-half assignments and count their residues modulo $M$.
  \item Enumerate all left-half assignments, respecting the leading-digit condition, and add the number of matching right-half residues.
  \item Sum over all lengths.
\end{enumerate}

\section{Complexity Analysis}

For a fixed length $d$, the larger half has size at most $8$, so the meet-in-the-middle phase costs
\[
  O(10^{h_1} + 10^{h_2}),
\]
with $\max(h_1,h_2)\le 8$. The memory usage is $O(M)$ for the residue-count table.

\section{Answer}

\[
\boxed{2000008332}
\]

\end{document}
